\documentclass[11pt,a4paper]{article}
\usepackage[margin=25mm]{geometry}
\usepackage[T1]{fontenc}
\usepackage{amsmath,amssymb,booktabs,longtable,array}
\usepackage[hidelinks]{hyperref}
\usepackage{xurl}
\newcommand{\code}[1]{\texttt{\nolinkurl{#1}}}
\newcommand{\ptmiss}{\boldsymbol{p}_{T}^{\mathrm{miss}}}
\newcommand{\res}{\boldsymbol{\delta}}
\newcommand{\sigq}{\sigma_{\mathrm{IQR}}}
\title{Implementation note: forward-cluster MET resolution study in $W\to\mu\nu$}
\author{}
\date{Code inspected on 9 September 2026}
\begin{document}
\maketitle

\section{Scope and conventions}
This note describes the implementation in
\code{analyses/met_resolution_study_wmunu.py},
\code{analyses/met_inputs_common.py}, and
\code{analyses/met_resolution_common.py}.
It documents the configured analysis, rather than prescribing a general MET
resolution method. The W entry point uses the truth W momentum as the boson
reference. Shared names such as \code{ptz}, \code{pz}, \code{pTZ}, and $C^Z$
are retained from the common W/Z implementation; they do not imply that a Z
is reconstructed in this study. In particular, \code{pz} denotes a transverse
projection, not momentum along the beam.

All momentum vectors below have two Cartesian components in GeV. Azimuthal
angles are used in radians. Input object transverse momenta, truth MET
transverse momenta, truth-boson transverse momentum, and track-soft MET
components are multiplied by $10^{-3}$. Cluster energies have the
scale-specific conversions listed in Section~\ref{sec:clusters}.
The study does not apply event weights in its retained distributions, means,
or quantiles. Stored hard-object weights are used in MET reconstruction;
these are distinct from event weights.

\section{Input processing and event selection}
Input files are supplied through repeated \code{--input} arguments and/or
recursive discovery of files ending in \code{.root} under \code{--input-dir}.
The combined paths are deduplicated and sorted. The default tree is
\code{treeAnaWZ}; \code{--nEvents} (also \code{--n-events}) limits processed
input events, with a negative default meaning all events. Files are processed
in chunks through \code{NtupleProcessor}.

For each event, the following operations occur in order:
\begin{enumerate}
\item Build the electron, muon, jet, and track-soft MET inputs. Require
consistent lengths of each hard-object term's $p_T$, $\phi$, and weight arrays,
including the jet $\eta$ array, and finite soft $x$ and $y$ components.
\item Require at least one entry in \code{mu_pt}. Despite the counter name
\code{one_muon}, this is not an exactly-one-muon requirement. No additional
muon kinematic, isolation, charge, or transverse-mass cut is implemented here.
\item Read \code{tboson_pt}, \code{tboson_eta}, and \code{tboson_phi}.
Reject nonfinite values and zero boson $p_T$, then require
$0\leq |\eta_W|<4.9$. The implementation tests $p_T\ne0$, assuming the physical
input convention of nonnegative transverse momentum.
\item Reconstruct the baseline MET and build a finite truth-MET reference.
Require that a truth transverse vector can be obtained.
\item Record the baseline truth-subtracted projection in a separate
``jet-inclusive'' collection, before cluster-array validation.
\item Identify forward input jets. Their presence increments a diagnostic
counter; there is no requirement that an event contain a forward jet.
\item Build the global cluster collection. Require the $\phi$ array and every
configured energy array to have the same length as the cluster $\eta$ array.
\item Store $|\eta_W|$, \code{avgMu}, and \code{nPrimVtx}, and fill all MET
variants for the event.
\end{enumerate}
Thus the ordinary variant histograms share the sample passing cluster-array
validation. The separate jet-inclusive mean can contain additional events.
There is no Z-mass selection in the W entry point.

\section{Baseline MET and truth reference}
For a hard-object term $a\in\{e,\mu,j\}$, define the visible vector
\begin{equation}
 \boldsymbol{H}_a=\sum_{i\in a}w_i p_{T,i}
       (\cos\phi_i,\sin\phi_i).
\end{equation}
The inputs are \code{met_input_el_pt/phi/weight},
\code{met_input_mu_pt/phi/weight}, and
\code{met_input_jet_pt/phi/weight}, with jet directions additionally taken
from \code{met_input_jet_eta}. The slash notation abbreviates separate branch
names with the same prefix. Objects with nonfinite $p_T$, $\phi$, or weight
are omitted from their vector sum. An empty valid sum is zero.

The baseline is reconstructed from these terms:
\begin{equation}
 \ptmiss_0=-(\boldsymbol{H}_e+\boldsymbol{H}_\mu+\boldsymbol{H}_j)
                 +\boldsymbol{S}_{\mathrm{trk}},
 \qquad
 \boldsymbol{S}_{\mathrm{trk}}=10^{-3}
 (\code{met_softtrk_mpx},\code{met_softtrk_mpy}).
\end{equation}
The stored soft term already has the missing-momentum sign. The study does
not use \code{met_maker_mpx/mpy} as its baseline, although a builder for those
branches exists in the common inputs file.

Truth MET is read from \code{tu_pt.second} and \code{tu_phi.second}.
The configured flattened-map indices are
$\mathrm{Int}:0$, $\mathrm{IntMuons}:1$, $\mathrm{IntOut}:2$, and
$\mathrm{NonInt}:3$. Currently
\code{TRUTH_MET_REFERENCE_COMPONENTS = ("NonInt",)}, so
\begin{equation}
 \ptmiss_{\mathrm{truth}}=10^{-3}p_{T,\mathrm{NonInt}}
    (\cos\phi_{\mathrm{NonInt}},\sin\phi_{\mathrm{NonInt}}).
\end{equation}
The builder supports summing several configured components, but does not do
so with the present setting. Some labels and docstrings still say
``Int + IntMuons''; the active configuration, rather than those labels,
determines the reference.

\section{Forward clusters and MET variants}\label{sec:clusters}
Forward input jets and selected forward clusters both satisfy
$2.5\leq|\eta|<4.9$, with finite $\eta$.
Cluster directions come from \code{fClusterEta} and \code{fClusterPhi}.
For each scale $s$, the code converts the energy to GeV as follows:
\begin{center}
\begin{tabular}{lll}
\toprule
Scale key & Energy branch & Multiplier to GeV\\
\midrule
\code{em} & \code{fCluster_rawE} & $10^{-3}$\\
\code{lcw} & \code{fCluster_calE} & $10^{-3}$\\
\code{ml} & \code{fCluster_MLE} & $10^{-3}$\\
\code{ml_forward} & \code{cluster_e_ML_forward} & $1$\\
\code{truth} & \code{fCluster_truthE} & $10^{-3}$\\
\bottomrule
\end{tabular}
\end{center}
For a selected cluster set $A$, its visible transverse vector is
\begin{equation}
 \boldsymbol{K}_s(A)=\sum_{k\in A}
 \frac{E_{k,s}}{\cosh\eta_k}(\cos\phi_k,\sin\phi_k).
\end{equation}
Only finite $\eta_k$, $\phi_k$, and $E_{k,s}$ enter this sum. Energy validity
is checked separately for each scale. There is no explicit energy threshold,
positivity cut, or additional cluster weight. Finite negative energies retain
their sign. The arrays are assumed to describe the same clusters in the same
order across scales; length checks alone cannot verify that assumption.

Let $F$ contain all forward clusters, and $J_F$ all forward input jets.
The non-overlapping subset $A\subset F$ satisfies
\begin{equation}
 \min_{j\in J_F}\Delta R_{kj}\geq0.4,
 \qquad \Delta R_{kj}^2=(\eta_k-\eta_j)^2+
 \left[\operatorname{atan2}(\sin(\phi_k-\phi_j),
                              \cos(\phi_k-\phi_j))\right]^2.
\end{equation}
If there are no forward jets, every forward cluster passes this overlap
requirement. Overlap is checked against forward input jets only, not against
central jets or leptons. These are geometrically selected global clusters,
not an independently tagged soft-cluster collection.

There are twelve variants: the baseline, a truth-MET control, and two
modifications for each of the five energy scales:
\begin{align}
 \ptmiss_{\mathrm{current}} &= \ptmiss_0,\\
 \ptmiss_{\mathrm{truth\ control}} &= \ptmiss_{\mathrm{truth}},\\
 \ptmiss_{\mathrm{keep},s} &= \ptmiss_0-\boldsymbol{K}_s(A),\\
 \ptmiss_{\mathrm{replace},s} &= \ptmiss_0+
                  \boldsymbol{H}_{j,F}-\boldsymbol{K}_s(F).
\end{align}
Here $\boldsymbol{H}_{j,F}$ uses the original forward jets' stored weights.
Adding it back removes their contribution to baseline MET. Subtracting a
cluster vector adds visible momentum to the reconstruction. The existing
track-soft term is retained in both modifications.
The variant keys are \code{current}, \code{truth_met},
\code{keep_jets_forward_clusters_<scale>}, and
\code{replace_forward_jets_with_all_forward_clusters_<scale>}.

\section{Residual and boson projection}
For each variant $v$, define
\begin{align}
 \res_v &= \ptmiss_v-\ptmiss_{\mathrm{truth}},\\
 \boldsymbol{q}_T &= 10^{-3}\code{tboson_pt}
                      (\cos\phi_W,\sin\phi_W),\qquad
 \widehat{\boldsymbol{q}}_T=\boldsymbol{q}_T/|\boldsymbol{q}_T|,\\
 P_{\parallel,v} &= \res_v\cdot\widehat{\boldsymbol{q}}_T
       =\delta_{x,v}\cos\phi_W+\delta_{y,v}\sin\phi_W.
\end{align}
All variants receive the same event's boson vector and truth reference.
The concrete call chain is
\begin{verbatim}
_append_variant(name, boson.vector, variant_met, truth_met_vector)
    residual = met - truth_met
    self.values[name].append(boson, met, residual)

ObservableValues.append(z_vector, met, projection_met)
    projected = met if projection_met is None else projection_met
    self.pz.append(projected.dot(z_vector.unit()))
\end{verbatim}
The third argument is therefore the residual, and truth subtraction occurs
before projection. The code also stores $|\ptmiss_v|$, its azimuth, $q_T$,
and the Cartesian residual components. The Cartesian $x$ residual is not
rotated into the boson frame.

\section{Binned response and resolution}
For a bin $B$ of the chosen variable, the response is
\begin{equation}\label{eq:response}
 C_v^Z(B)=1+\frac{\operatorname{median}(P_{\parallel,v}\mid B)}
                         {\langle q_T\rangle_B}.
\end{equation}
This is the implemented definition, including its plus sign and its ratio
of a median to a mean. It is not an average of event-by-event ratios.
The code computes it independently in boson-$p_T$ bins and in
\code{nPrimVtx} bins. It does not import a separately measured Z calibration.

The local quantile wrapper documents the robust width as
\begin{equation}\label{eq:width}
 \sigq(X\mid B)=\frac{Q_{84}(X\mid B)-Q_{16}(X\mid B)}{2}.
\end{equation}
This is approximately a Gaussian standard deviation. The code's name
``IQR'' refers to this central-68\% width, not the conventional
interquartile range $Q_{75}-Q_{25}$. Specifically, the wrapper calls
\code{calculate_sigma_iqr_in_x_bins} and divides its returned width by two.
The imported helper was outside the three-file inspection scope: its exact
percentile probabilities, interpolation, boundary handling, and minimum
population rules have not been independently verified in this note.

The response-corrected Cartesian resolution is
\begin{equation}
 R_{x,v}(B)=\frac{\sigq(\delta_{x,v}\mid B)}{C_v^Z(B)}.
\end{equation}
The division takes place after computing the width, using a single response
per bin. It does not first correct each event with a $p_T$-dependent response.
The numerator and response use the same binning variable and edges. Response
bins are initialized to NaN and filled only for finite median and finite,
nonzero mean $q_T$. Corrected widths require a finite numerator and finite,
nonzero response. No absolute value or positivity requirement is imposed on
the response, and no small-response threshold is applied.

The projected-width outputs named \code{raw_resolution} and
\code{resolution} both contain the same uncorrected
$\sigq(P_\parallel\mid B)$ in the current implementation.
The helper \code{calculate_binned_performance} also computes
\code{np.std}, using the population standard deviation, and a mean-based
response. These standard-deviation results are not used as the contents of
the MET-resolution histograms described here. The helper supplies counts,
and its arithmetic mean supplies the separate jet-inclusive mean histogram.

\section{Binning and output inventory}
The configured bin edges are
\begin{align*}
 q_T\ [\mathrm{GeV}]:\;&0,10,20,30,40,50,60,70,80,90,100,120,160,200,260,\\
 \code{avgMu}:\;&0,5,10,\ldots,80,\\
 \code{nPrimVtx}:\;&0,2,4,\ldots,50.
\end{align*}
The locally implemented binned mean and moment helper use half-open bins,
with the final upper edge included. Quantile-bin details are delegated to
the imported helper noted above. No $q_T<260$ event cut is applied before
the pileup calculations: vertex bins can include events beyond the plotted
boson-$p_T$ range.

In the following inventory, \code{<v>} denotes a variant key.
\begin{description}
\item[Basic distributions.]
\code{h_met_<v>} stores MET magnitude (120 bins, 0--300 GeV);
\code{h_met_phi_<v>} stores MET azimuth (64 bins, $-\pi$ to $\pi$);
\code{h_pz_<v>} stores $P_\parallel$ (120 bins, $-120$ to 120 GeV).
\code{h2_pz_vs_pTZ_<v>} uses the configured $q_T$ edges and the same
projection range. MET and projection distributions use a joint finite mask
on $q_T$, $P_\parallel$, and MET magnitude; the azimuth uses its own finite mask.
\item[Boson-$p_T$ summaries.]
\code{h_entries_vs_pTZ_<v>} contains event counts;
\code{h_mean_pz_vs_pTZ_<v>} contains the median projection despite its name;
\code{h_response_vs_pTZ_<v>} contains Eq.~\eqref{eq:response};
\code{h_raw_resolution_vs_pTZ_<v>} and
\code{h_resolution_vs_pTZ_<v>} both contain Eq.~\eqref{eq:width} for
$P_\parallel$.
\code{h_iqr_met_px_residual_over_cz_vs_pTZ_<v>} contains $R_x$.
\item[Pileup summaries.]
\code{h_rms_met_px_residual_vs_avgMu_<v>} and
\code{h_rms_met_py_residual_vs_avgMu_<v>} contain the uncorrected
robust widths of $\delta_x$ and $\delta_y$. Corresponding names with
\code{nPrimVtx} replace \code{avgMu}. The legacy \code{rms} names do not
indicate an RMS calculation.
\code{h_response_vs_nPrimVtx_<v>} stores the exact vertex-binned response
used by \code{h_iqr_met_px_residual_over_cz_vs_nPrimVtx_<v>}.
There is no analogous corrected $y$ or corrected \code{avgMu} output here.
\item[Shared distributions.]
\code{h_boson_eta} stores $|\eta_W|$ (100 bins, 0--4.9).
\code{h_mean_pz_vs_pTZ_jet_inclusive} stores the arithmetic mean of the
baseline projection before cluster-array validation. Additional forward
cluster diagnostics are delegated to \code{ForwardClusterDiagnostics};
their definitions lie outside this note's inspected source files.
\end{description}

The default output is \code{met_resolution_study.root}, configurable with
\code{--output}. It is opened in ROOT \code{RECREATE} mode. Observables are
accumulated, converted to NumPy arrays, reduced into histogram contents,
and written. The local analysis does not explicitly calculate uncertainties
on the quantiles or response ratios; ROOT histogram construction, including
any default errors and treatment of NaN contents, is delegated to imported
histogram helpers.

\section{Interpretation and implementation limits}
A response near unity means that the median projected residual is small
relative to the mean boson transverse momentum in that bin. It does not
imply a narrow Cartesian residual distribution. Different scales can have
similar $C^Z$ and different $R_x$.
Moreover, a response recomputed in vertex bins is not an average of responses
in boson-$p_T$ bins: the event mixtures differ, and the median-to-mean ratio
is nonlinear. The differing $q_T$ coverage noted above can also matter.

For finite, populated bins with nonzero mean boson $p_T$, the truth-MET
control has zero residual, unit response, and zero resolution. Deviations
from this expectation would warrant checking the input values and helper
behavior.
The scale conversions assume the documented branch units; those units and
cross-scale cluster ordering are not inferred or validated from the data.
Nonfinite cluster energies are skipped separately by scale, and overlap
removal only considers forward jets. These details should be retained when
comparing this implementation with other MET definitions.
\end{document}
